\documentclass[11pt,oneside]{article}

\usepackage[letterpaper,margin=0.88in,headheight=15pt]{geometry}
\usepackage[T1]{fontenc}
\usepackage[utf8]{inputenc}
\IfFileExists{libertinus.sty}{\usepackage{libertinus}}{\usepackage{lmodern}}
\usepackage{microtype}
\usepackage{amsmath,amssymb,amsthm,mathtools,mathrsfs,bm}
\usepackage{booktabs,longtable,tabularx,array,multirow}
\usepackage{enumitem}
\usepackage{xcolor}
\usepackage{graphicx}
\usepackage{float}
\usepackage{fancyhdr}
\usepackage{titlesec}
\usepackage[most]{tcolorbox}
\usepackage{listings}
\usepackage{xurl}
\usepackage[hypertexnames=false]{hyperref}
\usepackage[nameinlink,noabbrev]{cleveref}

\definecolor{linkblue}{RGB}{22,77,132}
\definecolor{deepgreen}{RGB}{28,103,76}
\definecolor{deepred}{RGB}{150,42,38}
\definecolor{warmgray}{RGB}{94,87,76}
\definecolor{paleblue}{RGB}{240,246,252}
\definecolor{palegreen}{RGB}{239,248,244}
\definecolor{paleyellow}{RGB}{254,250,233}
\definecolor{palered}{RGB}{253,241,240}
\definecolor{shadegray}{RGB}{246,247,249}
\definecolor{rulegray}{RGB}{190,197,204}

\hypersetup{
  colorlinks=true,
  linkcolor=linkblue,
  citecolor=deepgreen,
  urlcolor=linkblue,
  pdfauthor={Research report prepared for Vladimir Reshetnikov},
  pdftitle={Dyadic-Comb Sums for the Fabius--Rvachev System},
  pdfsubject={Exact monomial sums, Euler--Maclaurin collapse, complex powers, and iterated summation},
  pdfkeywords={Fabius function, Rvachev up-function, Thue--Morse, dyadic comb, Euler--Maclaurin, Hurwitz zeta, Bell polynomial, iterated sum}
}

\pagestyle{fancy}
\fancyhf{}
\fancyhead[L]{\small Dyadic-comb sums for Fabius--Rvachev functions}
\fancyhead[R]{\small\nouppercase{\leftmark}}
\fancyfoot[C]{\thepage}
\renewcommand{\headrulewidth}{0.4pt}

\titleformat{\section}{\Large\bfseries}{\thesection}{0.7em}{}
\titleformat{\subsection}{\large\bfseries}{\thesubsection}{0.7em}{}
\titleformat{\subsubsection}{\normalsize\bfseries}{\thesubsubsection}{0.7em}{}
\setcounter{secnumdepth}{3}
\setcounter{tocdepth}{2}
\setlist[itemize]{leftmargin=2em,itemsep=0.25em,topsep=0.4em}
\setlist[enumerate]{leftmargin=2.25em,itemsep=0.25em,topsep=0.4em}
\setlength{\emergencystretch}{3em}
\allowdisplaybreaks[3]
\numberwithin{equation}{section}

\newtheorem{theorem}{Theorem}[section]
\newtheorem{proposition}[theorem]{Proposition}
\newtheorem{lemma}[theorem]{Lemma}
\newtheorem{corollary}[theorem]{Corollary}
\newtheorem{identity}[theorem]{Identity}
\newtheorem{conjecture}[theorem]{Conjecture}
\theoremstyle{definition}
\newtheorem{definition}[theorem]{Definition}
\newtheorem{example}[theorem]{Example}
\newtheorem{algorithm}[theorem]{Algorithm}
\newtheorem{problem}[theorem]{Problem}
\theoremstyle{remark}
\newtheorem{remark}[theorem]{Remark}
\newtheorem{warning}[theorem]{Warning}
\crefname{theorem}{theorem}{theorems}
\Crefname{theorem}{Theorem}{Theorems}
\crefname{proposition}{proposition}{propositions}
\Crefname{proposition}{Proposition}{Propositions}
\crefname{lemma}{lemma}{lemmas}
\Crefname{lemma}{Lemma}{Lemmas}
\crefname{corollary}{corollary}{corollaries}
\Crefname{corollary}{Corollary}{Corollaries}
\crefname{identity}{identity}{identities}
\Crefname{identity}{Identity}{Identities}
\crefname{conjecture}{conjecture}{conjectures}
\Crefname{conjecture}{Conjecture}{Conjectures}
\crefname{definition}{definition}{definitions}
\Crefname{definition}{Definition}{Definitions}
\crefname{example}{example}{examples}
\Crefname{example}{Example}{Examples}
\crefname{algorithm}{algorithm}{algorithms}
\Crefname{algorithm}{Algorithm}{Algorithms}
\crefname{problem}{problem}{problems}
\Crefname{problem}{Problem}{Problems}
\crefname{remark}{remark}{remarks}
\Crefname{remark}{Remark}{Remarks}
\crefname{warning}{warning}{warnings}
\Crefname{warning}{Warning}{Warnings}
\newenvironment{proofidea}{\par\smallskip\noindent\textit{Proof idea.}\ }{\hfill$\square$\par\smallskip}

\newtcolorbox{resultbox}[1][]{
  enhanced,breakable,colback=paleblue,colframe=linkblue,
  boxrule=0.7pt,arc=1.5mm,left=1.5mm,right=1.5mm,top=1mm,bottom=1mm,
  title={#1},fonttitle=\bfseries
}
\newtcolorbox{statusbox}[1][]{
  enhanced,breakable,colback=palegreen,colframe=deepgreen,
  boxrule=0.7pt,arc=1.5mm,left=1.5mm,right=1.5mm,top=1mm,bottom=1mm,
  title={#1},fonttitle=\bfseries
}
\newtcolorbox{warningbox}[1][]{
  enhanced,breakable,colback=paleyellow,colframe=warmgray,
  boxrule=0.7pt,arc=1.5mm,left=1.5mm,right=1.5mm,top=1mm,bottom=1mm,
  title={#1},fonttitle=\bfseries
}
\newtcolorbox{conjecturebox}[1][]{
  enhanced,breakable,colback=palered,colframe=deepred,
  boxrule=0.7pt,arc=1.5mm,left=1.5mm,right=1.5mm,top=1mm,bottom=1mm,
  title={#1},fonttitle=\bfseries
}

\lstdefinestyle{python}{
  language=Python,
  basicstyle=\ttfamily\footnotesize,
  backgroundcolor=\color{shadegray},
  frame=single,
  rulecolor=\color{rulegray},
  columns=fullflexible,
  keepspaces=true,
  showstringspaces=false,
  breaklines=true,
  keywordstyle=\color{linkblue}\bfseries,
  commentstyle=\color{deepgreen},
  stringstyle=\color{deepred},
  xleftmargin=0.6em,xrightmargin=0.6em
}

\newcommand{\R}{\mathbb R}
\newcommand{\C}{\mathbb C}
\newcommand{\Q}{\mathbb Q}
\newcommand{\N}{\mathbb N}
\newcommand{\Z}{\mathbb Z}
\newcommand{\E}{\mathbb E}
\newcommand{\Prob}{\mathbb P}
\newcommand{\dd}{\,\mathrm d}
\newcommand{\e}{\mathrm e}
\newcommand{\ii}{\mathrm i}
\newcommand{\Up}{\operatorname{up}}
\newcommand{\sincpi}{\operatorname{sinc}_{\pi}}
\newcommand{\bigO}{\mathcal O}
\newcommand{\smallo}{o}
\newcommand{\eps}{\varepsilon}
\newcommand{\vtwo}{\nu_2}
\newcommand{\wt}{s_2}
\newcommand{\fall}[2]{#1^{\underline{#2}}}
\newcommand{\rise}[2]{#1^{\overline{#2}}}
\newcommand{\pospart}[1]{\left(#1\right)_{+}}
\newcommand{\repo}{\texttt{Analysis/FabiusFunction/docs}}
\newcommand{\Dop}{\mathcal D}
\newcommand{\Bell}{Y}
\newcommand{\Lcomb}{L}
\newcommand{\Tcomb}{T}
\newcommand{\Ucomb}{U}
\newcommand{\Pcomb}{P}
\newcommand{\Frow}{\mathscr F}
\newcommand{\Arow}{A}
\newcommand{\ThetaTM}{\Theta}
\newcommand{\defect}{\mathcal R}
\newcommand{\stirlingfirst}[2]{\genfrac{[}{]}{0pt}{}{#1}{#2}}
\DeclareMathOperator{\Res}{Res}
\DeclareMathOperator{\Li}{Li}
\DeclareMathOperator{\sgn}{sgn}

\title{\bfseries Dyadic-Comb Sums for the Fabius--Rvachev System\\[0.35em]
\Large Rational sample generating functions, exact Euler--Maclaurin collapse,\\
complex powers, and iterated summation}
\author{Research report prepared for Vladimir Reshetnikov}
\date{28 August 2026\\[0.25em]
\small Repository audit pinned to commit \texttt{fa53516af8cbe7d792a97b25378b228b96e37cf9}}

\begin{document}
\maketitle

\begin{abstract}
Let \(F\) be the Fabius distribution function, let \(u=\Up\) be Rvachev's even
probability density on \([-1,1]\), and put \(M=2^N\), \(h=M^{-1}\).  This paper
studies the unweighted dyadic-comb sums
\[
 L_{N,\alpha}=h\sum_{k=0}^{M-1}\left(\frac{k}{M}\right)^{\alpha}
 F\!\left(\frac{k}{M}\right),
\]
where the zero term is omitted for nonintegral or negative \(\alpha\).  These are ordinary
left Riemann sums, not the self-weighted up-function quadrature already present in the
repository corpus.

For every nonnegative integer \(p\), an exact finite Bernoulli--Thue--Morse formula is
obtained by inserting the dyadic evaluator for \(F(k/2^N)\) and performing the remaining
power sums by Faulhaber's formula.  A stronger compression follows from a Bell-polynomial
formula for complete signed Thue--Morse power blocks: although the original expression
contains \(2^N\) signs, only \(p+2\) Bell terms survive.  The entire sample row also has a
single rational ordinary generating function.  Its apparent pole of order \(N+1\) is
cancelled by the finite Thue--Morse product, leaving one simple pole whose constant tail is
exactly \(1\).

The central theorem is an exact Euler--Maclaurin collapse.  If
\(I_p=\int_0^1x^pF(x)\,\dd x\), then
\[
 L_{N,p}=I_p-\frac h2+
 \sum_{r=1}^{\lfloor(p+1)/2\rfloor}
 \frac{B_{2r}}{(2r)!}\fall p{2r-1}h^{2r}
\]
for every \(p\le 2\lfloor N/2\rfloor\).  Thus a normally asymptotic
Euler--Maclaurin series becomes a finite exact identity on a degree range that grows with
the dyadic level.  The proof uses periodic trapezoidal aliasing and the order
\(N+\nu_2(\ell)\) zeros of the Fabius characteristic function at frequencies
\(2^N\ell\).  At odd levels the first failure is an explicit half-integer sinc-product
Dirichlet series whose signs are Thue--Morse signs.

For arbitrary \(\alpha\in\C\), the finite comb is entire in \(\alpha\) and admits an
exact Hurwitz-zeta formula; negative integer powers are therefore exact rational harmonic
sums.  The corresponding Rvachev combs exhibit a critical logarithmic divergence at
\(\alpha=-1\) and zeta-dominated power divergences below it.  Finally, an \(n\)-fold
prefix sum is reduced exactly to finitely many shifted one-fold combs through unsigned
Stirling numbers of the first kind, and a fractional-order extension is given by the
binomial/Gamma summation kernel.  Exact rational and high-precision experiments accompany
the proofs.  Novelty claims are relative to the pinned repository snapshot and a targeted
literature search; no unconditional claim of worldwide priority is made.
\end{abstract}

\tableofcontents
\clearpage

\section{Scope, corpus audit, and contribution boundary}
\label{sec:scope}

\subsection{Pinned source boundary}

The requested source boundary was the recursive LaTeX tree under \repo{}.  The audit was
pinned to commit
\[
 \texttt{fa53516af8cbe7d792a97b25378b228b96e37cf9},
\]
dated 28 August 2026.  The repository's own source ledger records a stable 67-document
core together with live additions, while the pinned tree also contains auxiliary TeX
fragments and source-audit copies.  The audit therefore proceeded by mathematical
lineage rather than by treating byte-identical or embedded copies as independent prior
art.  In particular, the following consolidated families were used as the principal
interfaces:
\begin{itemize}
\item the primary Fabius--Rvachev exposition and mathematical glossary;
\item the dyadic, \(q\)-binomial, asymptotic, and Thue--Morse archive studies;
\item the consolidated inverse-and-sampling volume;
\item the consolidated integration-and-transform volume;
\item both representation-frontier volumes, including their embedded source audits;
\item the imported or corrected classical papers and the Lean walkthrough.
\end{itemize}

A targeted phrase and formula search was then made for ordinary dyadic Riemann sums,
unweighted Fabius combs, Euler--Maclaurin exactness, complete row generating functions,
and iterated comb sums.  No counterpart of the combined theorem package below was found.
The external search likewise located classical Fabius arithmetic and compact-support
foundations~\cite{Arias1982,AriasArithmetic}, general Euler--Maclaurin theory including
singular variants~\cite{BerlineVergne,BuchheitKessler}, and the standard Bernoulli--Hurwitz
special-function machinery~\cite{NIST}, but not the same construction.

\begin{statusbox}[Meaning of the status labels]
\textbf{Corpus-known} means explicitly present in the pinned repository and used here as
input.  \textbf{Proved here} means derived in this report from those inputs and standard
analysis.  \textbf{Experimentally checked} means verified by exact rational or
high-precision code but not promoted beyond the accompanying proof.  \textbf{Conjectural}
means that a plausible nonvanishing, sign, or asymptotic assertion remains open.
\end{statusbox}

\subsection{Strong negative controls: what is not claimed as new}

The repository already contains, in substantially greater generality than needed here,
the results summarized in its primary and consolidated volumes~\cite{ProveItPrimary,ProveItSampling,ProveItIntegration,ProveItRepresentations}:
\begin{itemize}
\item exact dyadic values of \(F\) through moments, Thue--Morse sums, and Gaussian
      \(q\)-binomial coefficients;
\item moment recurrences, Bernoulli cumulants, complete Bell-polynomial formulas, and
      inverse-dyadic arithmetic;
\item the Fourier product for \(u\), integer zero multiplicity \(1+\nu_2(m)\), local
      zero jets, half-integer signs, and sharp Fourier-decay analyses;
\item positive self-weighted shifted quadrature
      \(2^{-N}\sum p((k+\theta)/2^N)u((k+\theta)/2^N)\), exact for
      \(\deg p\le N\);
\item continuous integrals \(\int_0^1x^\alpha F(x)\,\dd x\), their Mellin continuation,
      negative powers, Volterra ladders, and fractional integrals;
\item repeated integration of Thue--Morse signs, geometric Richardson filters, and
      inverse-Fabius sampling identities.
\end{itemize}

The object studied here is different: the lattice is uniform and the summand itself is
\(x^\alpha F(x)\) or \(x^\alpha u(x)\); no sample of \(u\) is inserted as an
integration weight.  The exact Euler--Maclaurin collapse in \cref{sec:em} is therefore
not a restatement of the corpus's self-sampling quadrature theorem, although
\cref{sec:up} shows that the two theories meet in a useful consistency identity.

\subsection{Contribution map}

\begin{table}[H]
\centering
\small
\begin{tabularx}{\textwidth}{@{}p{0.27\textwidth}X p{0.18\textwidth}@{}}
\toprule
Topic & Result & Status \\
\midrule
Sample-row algebra & Rational OGF for the complete level-\(N\) Fabius row, with exact
pole cancellation and constant tail. & Proved here \\
Integer powers & Finite Bernoulli--Thue--Morse formula for every \(N,p\), plus a
\(p+2\)-term Bell--Prouhet compression. & Proved here \\
Riemann-sum exactness & Euler--Maclaurin is exact for
\(p\le2\lfloor N/2\rfloor\). & Proved here \\
First failure & Odd-level leading defect as a half-integer sinc/Thue--Morse Dirichlet
series. & Proved here \\
Complex powers & Entire finite comb and exact Hurwitz-zeta representation; rational
negative-integer specializations. & Proved here \\
Rvachev combs & Meromorphic continuous limit, critical logarithm at \(-1\), and
zeta-dominated lower powers. & Proved here \\
Iterated sums & Exact Stirling reduction of an \(n\)-fold sum to shifted one-fold combs;
Gamma-kernel fractional order. & Proved here \\
Exact precision & Nonvanishing of the first degree outside the proved range for every
level. & Conjectural \\
\bottomrule
\end{tabularx}
\caption{Main outputs and proof status.}
\label{tab:contribution-map}
\end{table}

\section{Normalization and inherited exact data}
\label{sec:background}

\subsection{The Fabius law and Rvachev density}

Let \(Y\) be the random variable
\begin{equation}
 Y=\sum_{j=1}^{\infty}2^{-j}U_j,
 \qquad U_j\stackrel{\mathrm{iid}}{\sim}\operatorname{Unif}[0,1],
 \label{eq:Y-law}
\end{equation}
and let
\[
 F(x)=\Prob(Y\le x),\qquad 0\le x\le1.
\]
Then
\begin{equation}
 F(1-x)=1-F(x),
 \qquad F'(x)=2F(2x)\quad(0\le x\le\tfrac12),
 \label{eq:F-basic}
\end{equation}
and \(F\) is flat at both endpoints.  Put \(X=2Y-1\).  Its density is Rvachev's
even up-function \(u=\Up\), normalized by
\begin{equation}
 \int_{-1}^{1}u(x)\,\dd x=1,
 \qquad
 u(x)=F(1-|x|),
 \qquad
 u(x)=1-F(x)\quad(0\le x\le1).
 \label{eq:up-fold}
\end{equation}

The Fourier convention used throughout is
\begin{equation}
 \Phi(\xi)=\int_{-1}^{1}u(x)\e^{-2\pi\ii x\xi}\,\dd x
 =\prod_{j=0}^{\infty}\sincpi\!\left(\frac{\xi}{2^j}\right),
 \qquad \sincpi z=\frac{\sin(\pi z)}{\pi z}.
 \label{eq:Phi-product}
\end{equation}
At every nonzero integer \(m\),
\begin{equation}
 \operatorname{ord}_{\xi=m}\Phi=1+\vtwo(m).
 \label{eq:Phi-zero-order}
\end{equation}
The characteristic function of \(Y\) is therefore
\begin{equation}
 \chi(\xi)=\E\e^{-2\pi\ii\xi Y}
 =\e^{-\pi\ii\xi}\Phi(\xi/2).
 \label{eq:chi-Phi}
\end{equation}

\subsection{Moments and continuous monomial integrals}

Write
\begin{equation}
 c_r=\int_{-1}^{1}x^{2r}u(x)\,\dd x,
 \qquad c_0=1.
 \label{eq:cr}
\end{equation}
The exact recurrence is
\begin{equation}
 (2r+1)(2^{2r}-1)c_r
 =\sum_{j=0}^{r-1}\binom{2r+1}{2j}c_j.
 \label{eq:c-recurrence}
\end{equation}
By \(Y=(X+1)/2\), the complete Fabius moments are
\begin{equation}
 d_m:=\E Y^m
 =2^{-m}\sum_{j=0}^{\lfloor m/2\rfloor}\binom m{2j}c_j.
 \label{eq:dm-from-c}
\end{equation}
For \(p\in\N_0\), integration by parts gives
\begin{equation}
 I_p:=\int_0^1x^pF(x)\,\dd x
 =\frac{1-d_{p+1}}{p+1}
 =\frac1{p+1}\left(
 1-2^{-p-1}\sum_{j=0}^{\lfloor(p+1)/2\rfloor}
 \binom{p+1}{2j}c_j\right).
 \label{eq:Ip}
\end{equation}
More generally, if
\begin{equation}
 d(s)=\E Y^s,
 \label{eq:entire-moment}
\end{equation}
then endpoint flatness makes \(d(s)\) entire, and
\begin{equation}
 I(\alpha):=\int_0^1x^\alpha F(x)\,\dd x
 =\frac{1-d(\alpha+1)}{\alpha+1}
 \label{eq:Ialpha}
\end{equation}
is entire after removing the apparent singularity at \(\alpha=-1\).

\subsection{The dyadic value interface}

Let
\begin{equation}
 M=2^N,
 \qquad h=M^{-1},
 \qquad \varepsilon_a=(-1)^{\wt(a)},
 \label{eq:M-h-eps}
\end{equation}
and define
\begin{equation}
 P_N(y)=\sum_{j=0}^{\lfloor N/2\rfloor}
 \binom N{2j}c_jy^{N-2j},
 \qquad
 C_N=\frac{2^{-N(N+1)/2}}{N!}.
 \label{eq:PN-CN}
\end{equation}
The corpus's exact dyadic evaluator can be written compactly as
\begin{equation}
 \boxed{\displaystyle
 F\!\left(\frac{k}{M}\right)
 =C_N\sum_{a=0}^{k-1}\varepsilon_a
 P_N\bigl(2(k-a)-1\bigr),
 \qquad 0\le k\le M.}
 \label{eq:dyadic-convolution}
\end{equation}
The empty sum at \(k=0\) is zero.  This finite convolution is the sole arithmetic
input needed for \cref{sec:row,sec:integer}.

\section{The dyadic combs}
\label{sec:def}

\subsection{Left, right, and trapezoidal forms}

For \(\alpha\in\C\), define the punctured left Fabius comb by
\begin{equation}
 \boxed{\displaystyle
 L_{N,\alpha}
 =\frac1M\sum_{k=1}^{M-1}
 \left(\frac{k}{M}\right)^\alpha F\!\left(\frac{k}{M}\right),}
 \label{eq:L-def}
\end{equation}
where positive real powers use the real logarithm.  For \(p\in\N_0\), this agrees with
the ordinary left Riemann sum including \(k=0\), because \(F(0)=0\).  At fixed
\(N\), \(L_{N,\alpha}\) is an exponential polynomial and hence an entire function of
\(\alpha\).

The right and trapezoidal sums for \(x^pF(x)\) are
\begin{equation}
 R_{N,p}=L_{N,p}+h,
 \qquad
 T_{N,p}=L_{N,p}+\frac h2,
 \label{eq:right-trap}
\end{equation}
because \(F(0)=0\) and \(F(1)=1\).  The trapezoidal normalization is the natural one
for Fourier aliasing.

\subsection{Shifted phases}

For a phase \(\theta\in\R\), one may similarly define
\begin{equation}
 L_{N,\alpha}^{(\theta)}
 =h\sum_{\substack{k\in\Z\\0<k+\theta<M}}
 \left(\frac{k+\theta}{M}\right)^\alpha
 F\!\left(\frac{k+\theta}{M}\right).
 \label{eq:shifted-comb}
\end{equation}
The present report emphasizes the centered phase \(\theta=0\), where all samples are
level-\(N\) dyadic values.  A dyadic phase \(\theta=a/2^s\) reduces to a residue class
inside the finer level \(N+s\); an exact roots-of-unity filter is given in
\cref{cor:roots-of-unity}.

\section{A rational generating function for a complete sample row}
\label{sec:row}

\subsection{Euler operator form}

Let
\begin{equation}
 \Dop=z\frac{\dd}{\dd z},
 \qquad
 \ThetaTM_N(z)=\prod_{j=0}^{N-1}(1-z^{2^j})
 =\sum_{a=0}^{M-1}\varepsilon_a z^a.
 \label{eq:thetaN}
\end{equation}
The kernel in \eqref{eq:dyadic-convolution} has OGF
\begin{equation}
 \sum_{r\ge1}P_N(2r-1)z^r
 =P_N(2\Dop-1)\frac{z}{1-z}.
 \label{eq:kernel-operator}
\end{equation}
Convolution with the Thue--Morse block therefore gives the following result.

\begin{theorem}[Complete sample-row generating function]
\label{thm:row-ogf}
Define
\begin{equation}
 \Frow_N(z)
 :=C_N\ThetaTM_N(z)P_N(2\Dop-1)\frac{z}{1-z}.
 \label{eq:Frow}
\end{equation}
Then
\begin{equation}
 \Frow_N(z)=\sum_{k\ge0}\widetilde F_N(k)z^k,
 \qquad
 \widetilde F_N(k)=
 \begin{cases}
 F(k/M),&0\le k\le M,\\
 1,&k\ge M.
 \end{cases}
 \label{eq:Frow-coeff}
\end{equation}
In particular,
\begin{equation}
 \boxed{\displaystyle
 A_N(z):=\sum_{k=0}^{M-1}F(k/M)z^k
 =\Frow_N(z)-\frac{z^M}{1-z}}
 \label{eq:AN}
\end{equation}
is a polynomial of degree at most \(M-1\).
\end{theorem}

\begin{proof}
The coefficient identity through \(k=M\) is exactly
\eqref{eq:dyadic-convolution}.  Applying \((2\Dop-1)^r\) to \(z/(1-z)\) produces a
rational function with denominator \((1-z)^{r+1}\) and numerator degree at most
\(r+1\).  Hence the largest denominator in
\(P_N(2\Dop-1)z/(1-z)\) is \((1-z)^{N+1}\).
On the other hand,
\begin{equation}
 \ThetaTM_N(z)=(1-z)^N
 \prod_{j=0}^{N-1}(1+z+\cdots+z^{2^j-1}).
 \label{eq:theta-factor}
\end{equation}
Thus \(\Frow_N\) has at most one simple pole at \(z=1\), and its numerator has degree
at most
\[
 (M-1)+(N+1)-N=M.
\]
Consequently its coefficients are constant from index \(M\) onward.  The coefficient
at index \(M\) is \(F(1)=1\), so the constant tail is \(1\).  Subtracting
\(z^M/(1-z)\) removes it and leaves \(A_N\).
\end{proof}

\begin{resultbox}[Pole collapse]
The finite Thue--Morse product does more than enforce Prouhet cancellation.  It removes
\emph{exactly} \(N\) of the \(N+1\) poles generated by the degree-\(N\) polynomial
kernel.  The surviving simple pole encodes the saturation \(F(k/M)=1\) after the right
endpoint.
\end{resultbox}

\begin{corollary}[Natural powers by Euler differentiation]
\label{cor:euler-natural}
For every \(p\in\N_0\),
\begin{equation}
 \boxed{\displaystyle
 L_{N,p}=M^{-p-1}\left.(\Dop^pA_N)(z)\right|_{z=1}.}
 \label{eq:euler-natural}
\end{equation}
If the spectral Euler power is defined on polynomials by
\begin{equation}
 \Dop^\alpha\!\left(\sum_{k\ge0}a_kz^k\right)
 :=\sum_{k\ge1}k^\alpha a_kz^k,
 \label{eq:fractional-Euler}
\end{equation}
then the same formula holds for every \(\alpha\in\C\):
\begin{equation}
 L_{N,\alpha}=M^{-\alpha-1}(\Dop^\alpha A_N)(1).
 \label{eq:euler-complex}
\end{equation}
\end{corollary}

\begin{corollary}[Dyadic phase as a roots-of-unity filter]
\label{cor:roots-of-unity}
Let \(\theta=a/2^s\), put \(D=N+s\), and let
\(\omega=\exp(2\pi\ii/2^s)\).  Then
\begin{equation}
 \boxed{\displaystyle
 L_{N,\alpha}^{(a/2^s)}
 =\frac{2^{-D\alpha}}{M2^s}
 \sum_{r=0}^{2^s-1}\omega^{-ar}
 (\Dop^\alpha A_D)(\omega^r),}
 \label{eq:root-filter}
\end{equation}
with the obvious endpoint convention.
\end{corollary}

\begin{proof}
The usual character projector
\(2^{-s}\sum_r\omega^{r(k-a)}\) selects indices \(k\equiv a\pmod{2^s}\) from the
finer level-\(D\) row.
\end{proof}

\subsection{Low-level examples}

The exact rational functions begin
\begin{align}
 \Frow_1(z)&=\frac{z(z+1)}{2(1-z)},
 \label{eq:row-N1}\\
 \Frow_2(z)&=\frac{z(z+1)(z+5)(5z+1)}{72(1-z)},
 \label{eq:row-N2}\\
 \Frow_3(z)&=\frac{z(z+1)^3(z^2+1)(z^2+16z+1)}{288(1-z)}.
 \label{eq:row-N3}
\end{align}
The numerator is reciprocal after removal of the initial \(z\), reflecting
\(F(k/M)+F((M-k)/M)=1\).  A systematic factorization theory for these row numerators is
proposed in \cref{sec:frontiers}.

\section{Closed forms for nonnegative integer powers}
\label{sec:integer}

\subsection{Bernoulli--Thue--Morse formula}

Insert \eqref{eq:dyadic-convolution} into \eqref{eq:L-def}, exchange the finite sums,
and put \(r=k-a\).  Expanding \((a+r)^p\) and \((2r-1)^d\) leaves ordinary power sums
\(\sum_{r=1}^{M-1-a}r^q\), which are Bernoulli-polynomial differences.

\begin{theorem}[Exact closed form for every \(N,p\)]
\label{thm:bernoulli-TM}
For \(p\in\N_0\),
\begin{equation}
 \boxed{\displaystyle
 L_{N,p}=C_NM^{-p-1}
 \sum_{a=0}^{M-1}\varepsilon_a H_{N,p}(a),}
 \label{eq:L-H}
\end{equation}
where
\begin{align}
 H_{N,p}(a)
 &=\sum_{j=0}^{\lfloor N/2\rfloor}\binom N{2j}c_j
   \sum_{r=0}^{p}\binom pr a^{p-r}
   \sum_{s=0}^{N-2j}\binom{N-2j}{s}2^s(-1)^{N-2j-s}
 \notag\\[-0.25em]
 &\qquad\qquad\qquad\times
 \frac{B_{r+s+1}(M-a)-B_{r+s+1}(1)}{r+s+1}.
 \label{eq:Hnp}
\end{align}
Here \(B_n(x)\) denotes the Bernoulli polynomial.  The term \(a=M-1\) vanishes
automatically.
\end{theorem}

\begin{proof}
From \eqref{eq:dyadic-convolution},
\[
 M^{p+1}C_N^{-1}L_{N,p}
 =\sum_{a=0}^{M-2}\varepsilon_a
 \sum_{r=1}^{M-1-a}(a+r)^pP_N(2r-1).
\]
Use
\[
 (a+r)^p=\sum_{q=0}^{p}\binom pq a^{p-q}r^q,
 \qquad
 (2r-1)^d=\sum_{s=0}^{d}\binom ds2^s(-1)^{d-s}r^s,
\]
and Faulhaber's identity
\begin{equation}
 \sum_{r=1}^{K}r^m
 =\frac{B_{m+1}(K+1)-B_{m+1}(1)}{m+1}.
 \label{eq:faulhaber-B1}
\end{equation}
Taking \(K=M-1-a\) gives \eqref{eq:Hnp}.
\end{proof}

\begin{remark}
Formula \eqref{eq:L-H} is manifestly rational.  It is useful for exact evaluation at
moderate \(N\), but it still contains a complete \(2^N\)-term signed block.  The next
subsection removes that exponential-looking part.
\end{remark}

\subsection{Bell--Prouhet collapse}

Define the signed block moments
\begin{equation}
 T_{N,m}=\sum_{a=0}^{M-1}\varepsilon_a a^m.
 \label{eq:TNm}
\end{equation}
Their exponential generating function is
\begin{equation}
 \sum_{m\ge0}T_{N,m}\frac{t^m}{m!}
 =\ThetaTM_N(\e^t)
 =\prod_{j=0}^{N-1}(1-\e^{2^jt}).
 \label{eq:T-egf}
\end{equation}
Factor each term as
\[
 1-\e^x=-x\exp(x/2)\frac{\sinh(x/2)}{x/2}.
\]
This gives a closed logarithmic jet.

\begin{theorem}[Bell formula for complete Thue--Morse power blocks]
\label{thm:Bell-Prouhet}
Let
\begin{align}
 \lambda_{N,1}&=\frac{M-1}{2},
 \label{eq:lambda1}\\
 \lambda_{N,2r}&=\frac{B_{2r}}{2r}
 \frac{2^{2rN}-1}{2^{2r}-1}\qquad(r\ge1),
 \label{eq:lambda-even}\\
 \lambda_{N,2r+1}&=0\qquad(r\ge1).
 \label{eq:lambda-odd}
\end{align}
Then
\begin{equation}
 \ThetaTM_N(\e^t)
 =(-1)^N2^{N(N-1)/2}t^N
 \exp\!\left(\sum_{r\ge1}\lambda_{N,r}\frac{t^r}{r!}\right).
 \label{eq:theta-Bell-egf}
\end{equation}
Consequently,
\begin{equation}
 T_{N,m}=0\quad(m<N),
 \label{eq:Prouhet-zero}
\end{equation}
and, for \(s\ge0\),
\begin{equation}
 \boxed{\displaystyle
 T_{N,N+s}=(-1)^N2^{N(N-1)/2}
 \frac{(N+s)!}{s!}
 \Bell_s(\lambda_{N,1},\ldots,\lambda_{N,s}),}
 \label{eq:TN-Bell}
\end{equation}
where \(\Bell_s\) is the complete exponential Bell polynomial.
\end{theorem}

\begin{proof}
The standard expansion
\begin{equation}
 \log\frac{\e^x-1}{x}
 =\frac x2+\sum_{r\ge1}\frac{B_{2r}}{2r}\frac{x^{2r}}{(2r)!}
 \label{eq:log-expm1}
\end{equation}
may be summed over \(x=2^jt\).  The leading factors contribute
\((-1)^N2^{0+1+\cdots+(N-1)}t^N\), the linear terms give
\((M-1)t/2\), and the even terms give the finite geometric sum in
\eqref{eq:lambda-even}.  Comparing with
\[
 \exp\!\left(\sum_{r\ge1}x_r\frac{t^r}{r!}\right)
 =\sum_{s\ge0}\Bell_s(x_1,\ldots,x_s)\frac{t^s}{s!}
\]
proves the coefficient formula.
\end{proof}

\begin{corollary}[Only \(p+2\) signed moments survive]
\label{cor:pplus2}
The polynomial \(H_{N,p}(a)\) has degree at most \(N+p+1\).  Therefore
\begin{equation}
 \boxed{\displaystyle
 L_{N,p}=C_NM^{-p-1}
 \sum_{s=0}^{p+1}[a^{N+s}]H_{N,p}(a)\,T_{N,N+s},}
 \label{eq:Bell-collapse}
\end{equation}
with \(T_{N,N+s}\) given by \eqref{eq:TN-Bell}.  Thus the original full
Thue--Morse block collapses to at most \(p+2\) Bell-polynomial terms, independently of
\(N\).
\end{corollary}

\begin{proof}
In \eqref{eq:Hnp}, the Bernoulli polynomial has degree \(r+s+1\); after multiplication
by \(a^{p-r}\), the degree is at most \(p+s+1\le p+N+1\).  All monomials below degree
\(N\) vanish by \eqref{eq:Prouhet-zero}.
\end{proof}

\begin{resultbox}[A useful computational dichotomy]
For fixed \(N\) and many degrees \(p\), the row polynomial \(A_N\) is convenient.  For
fixed \(p\) and large symbolic \(N\), the Bell collapse is preferable: its signed part
has constant length \(p+2\), and all dependence on the exponentially long block is
compressed into the explicit jets \(\lambda_{N,r}\).
\end{resultbox}

\section{Exact Euler--Maclaurin collapse}
\label{sec:em}

\subsection{Periodic trapezoidal aliasing}

Put
\begin{equation}
 g_p(x)=x^pF(x),\qquad 0\le x\le1,
 \label{eq:gp}
\end{equation}
and let
\begin{equation}
 \widehat g_p(n)=\int_0^1g_p(x)\e^{-2\pi\ii nx}\,\dd x.
 \label{eq:gphat}
\end{equation}
The periodic extension of \(g_p\) has endpoint jump \(1\).  Its Fourier series takes
the midpoint value at the jump, so the trapezoidal alias identity is
\begin{equation}
 T_{N,p}=\sum_{\ell\in\Z}\widehat g_p(M\ell).
 \label{eq:trap-alias}
\end{equation}
The zero-frequency term is \(I_p\).

A probability-layer computation makes every nonzero Fourier coefficient explicit.

\begin{lemma}[Fourier coefficient through the Fabius characteristic function]
\label{lem:gp-hat}
For \(n\in\Z\setminus\{0\}\),
\begin{equation}
 \boxed{\displaystyle
 \widehat g_p(n)
 =\sum_{j=0}^{p}
 \frac{\fall p j}{(2\pi\ii n)^{j+1}}
 \left(
 \frac{\chi^{(p-j)}(n)}{(-2\pi\ii)^{p-j}}-1
 \right).}
 \label{eq:gp-hat-chi}
\end{equation}
\end{lemma}

\begin{proof}
Since \(F(x)=\Prob(Y\le x)\), Fubini gives
\begin{equation}
 \widehat g_p(n)=\E\int_Y^1x^p\e^{-2\pi\ii nx}\,\dd x.
 \label{eq:layer-fourier}
\end{equation}
Repeated integration by parts, using \(\e^{-2\pi\ii n}=1\), yields
\[
 \int_Y^1x^p\e^{-2\pi\ii nx}\,\dd x
 =\sum_{j=0}^{p}\frac{\fall p j}{(2\pi\ii n)^{j+1}}
 \left(Y^{p-j}\e^{-2\pi\ii nY}-1\right).
\]
Finally,
\(
 \E[Y^r\e^{-2\pi\ii nY}]=\chi^{(r)}(n)/(-2\pi\ii)^r
\).
\end{proof}

\subsection{The exact remainder}

The \(-1\) terms in \eqref{eq:gp-hat-chi} can be summed by Euler's formula for
\(\zeta(2r)\); they are precisely the ordinary Euler--Maclaurin boundary corrections.
The remaining characteristic-function terms form an exact spectral remainder.

\begin{theorem}[Exact Euler--Maclaurin plus an alias remainder]
\label{thm:EM-remainder}
For every \(N\ge0\) and \(p\in\N_0\),
\begin{equation}
 \boxed{\displaystyle
 L_{N,p}=I_p-\frac h2+
 \sum_{r=1}^{\lfloor(p+1)/2\rfloor}
 \frac{B_{2r}}{(2r)!}\fall p{2r-1}h^{2r}
 +\defect_{N,p},}
 \label{eq:EM-exact-remainder}
\end{equation}
where
\begin{equation}
 \boxed{\displaystyle
 \defect_{N,p}
 =\sum_{\ell\in\Z\setminus\{0\}}
 \sum_{j=0}^{p}
 \frac{\fall p j}{(2\pi\ii M\ell)^{j+1}}
 \frac{\chi^{(p-j)}(M\ell)}{(-2\pi\ii)^{p-j}}.}
 \label{eq:R-alias}
\end{equation}
The series is absolutely convergent after pairing and, equivalently, is the Fourier
series remainder of the periodic trapezoidal rule.
\end{theorem}

\begin{proof}
Insert \eqref{eq:gp-hat-chi} into \eqref{eq:trap-alias}.  The characteristic terms give
\eqref{eq:R-alias}.  For the boundary terms, only odd \(j=2r-1\) survive the symmetric
sum over \(\ell\), and
\begin{align*}
 -\frac{\fall p{2r-1}}{(2\pi\ii M)^{2r}}
 \sum_{\ell\ne0}\ell^{-2r}
 &=\frac{B_{2r}}{(2r)!}\fall p{2r-1}M^{-2r}.
\end{align*}
Finally \(T_{N,p}=L_{N,p}+h/2\).
\end{proof}

\subsection{Growing exactness range}

At the dual lattice, \eqref{eq:chi-Phi} and \eqref{eq:Phi-zero-order} give, for
\(N\ge1\),
\begin{equation}
 \operatorname{ord}_{\xi=M\ell}\chi
 =N+\vtwo(\ell),
 \qquad \ell\ne0.
 \label{eq:chi-zero-order}
\end{equation}
This forces almost every term in \eqref{eq:R-alias} to vanish.

\begin{theorem}[Exact Euler--Maclaurin collapse]
\label{thm:EM-collapse}
Let
\begin{equation}
 d_N=2\left\lfloor\frac N2\right\rfloor.
 \label{eq:dN}
\end{equation}
Then, for every \(p\in\N_0\) with \(p\le d_N\),
\begin{equation}
 \boxed{\displaystyle
 L_{N,p}=I_p-\frac h2+
 \sum_{r=1}^{\lfloor(p+1)/2\rfloor}
 \frac{B_{2r}}{(2r)!}\fall p{2r-1}h^{2r}.}
 \label{eq:EM-collapse}
\end{equation}
Equivalently, the level-\(N\) dyadic comb agrees exactly with its complete finite
Euler--Maclaurin polynomial through degree \(N\) when \(N\) is even, and through degree
\(N-1\) when \(N\) is odd.
\end{theorem}

\begin{proof}
The case \(N=0,p=0\) follows from symmetry:
\(L_{0,0}=0=I_0-1/2\).  Suppose \(N\ge1\).  If \(p<N\), every derivative order
\(p-j\) occurring in \eqref{eq:R-alias} is below the zero order
\(N+\nu_2(\ell)\), so \(\defect_{N,p}=0\).

It remains to consider \(p=N\).  Aliases with even \(\ell\) still vanish.  At odd
\(\ell\), only the \(j=0\) term can survive.  Since \(\Phi\) is even,
\[
 \chi^{(N)}(-M\ell)=(-1)^N\chi^{(N)}(M\ell).
\]
The denominator in the \(j=0\) term changes sign under \(\ell\mapsto-\ell\).  Hence
positive and negative odd aliases cancel exactly when \(N\) is even.  For odd \(N\), no
such cancellation is forced.  This is precisely the range \(p\le2\lfloor N/2\rfloor\).
\end{proof}

\begin{remark}[Why this is not ordinary polynomial exactness]
The summand contains the nonpolynomial function \(F\).  The theorem does not say that a
fixed quadrature rule integrates all polynomials against a fixed weight.  Instead, the
boundary polynomial generated by Euler--Maclaurin becomes exact because the nonzero
Fourier aliases are killed by increasingly high-order zeros of \(\chi\).  It is a
spectral cancellation theorem for a family of ordinary Riemann sums.
\end{remark}

\begin{corollary}[Low-degree formulas]\label{cor:low-degree}
With \(h=2^{-N}\), the formulas listed in \cref{tab:integer-closed} are exact once the displayed level threshold is met.
\end{corollary}

\begin{table}[H]
\centering
\small
\resizebox{\textwidth}{!}{% Generated by dyadic_comb_experiments.py
\begin{tabular}{c c c l}
\toprule
$p$ & $I_p$ & sufficient $N$ & exact $L_{N,p}$ \\
\midrule
0 & $\frac{1}{2}$ & $0$ & $\frac{1}{2}-\frac{h}{2}$ \\
1 & $\frac{13}{36}$ & $2$ & $\frac{13}{36}-\frac{h}{2} + \frac{1}{12}h^{2}$ \\
2 & $\frac{5}{18}$ & $2$ & $\frac{5}{18}-\frac{h}{2} + \frac{1}{6}h^{2}$ \\
3 & $\frac{1207}{5400}$ & $4$ & $\frac{1207}{5400}-\frac{h}{2} + \frac{1}{4}h^{2} - \frac{1}{120}h^{4}$ \\
4 & $\frac{251}{1350}$ & $4$ & $\frac{251}{1350}-\frac{h}{2} + \frac{1}{3}h^{2} - \frac{1}{30}h^{4}$ \\
5 & $\frac{22661}{142884}$ & $6$ & $\frac{22661}{142884}-\frac{h}{2} + \frac{5}{12}h^{2} - \frac{1}{12}h^{4} + \frac{1}{252}h^{6}$ \\
6 & $\frac{16427}{119070}$ & $6$ & $\frac{16427}{119070}-\frac{h}{2} + \frac{1}{2}h^{2} - \frac{1}{6}h^{4} + \frac{1}{42}h^{6}$ \\
7 & $\frac{63446897}{520506000}$ & $8$ & $\frac{63446897}{520506000}-\frac{h}{2} + \frac{7}{12}h^{2} - \frac{7}{24}h^{4} + \frac{1}{12}h^{6} - \frac{1}{240}h^{8}$ \\
8 & $\frac{7096441}{65063250}$ & $8$ & $\frac{7096441}{65063250}-\frac{h}{2} + \frac{2}{3}h^{2} - \frac{7}{15}h^{4} + \frac{2}{9}h^{6} - \frac{1}{30}h^{8}$ \\
\bottomrule
\end{tabular}
}
\caption{Exact Euler--Maclaurin formulas generated and checked by the companion script.}
\label{tab:integer-closed}
\end{table}

\subsection{The first odd-level defect}

At odd \(N\), the first degree outside the exact range is \(p=N\), and only odd aliases
contribute.  The local factorization of \(\Phi\) at
\(m=2^{N-1}q\), \(q\) odd, gives
\begin{equation}
 \Phi^{(N)}(2^{N-1}q)
 =-\frac{N!}{(2^{N-1}q)^N}\Phi(q/2).
 \label{eq:leading-zero-jet}
\end{equation}

\begin{theorem}[First odd-level defect]
\label{thm:odd-defect}
If \(N\ge1\) is odd, then
\begin{equation}
 \boxed{\displaystyle
 \defect_{N,N}
 =\frac{2(-1)^{(N+1)/2}N!}
 {(2\pi)^{N+1}2^{N(N+1)}}
 \sum_{m=0}^{\infty}
 \frac{\Phi((2m+1)/2)}{(2m+1)^{N+1}}.}
 \label{eq:odd-defect-series}
\end{equation}
Moreover,
\begin{equation}
 \sgn\Phi((2m+1)/2)=(-1)^{\wt(m)}=\varepsilon_m,
 \label{eq:half-sign}
\end{equation}
so the defect is a Thue--Morse-signed half-integer sinc-product Dirichlet series.
\end{theorem}

\begin{proof}
For \(p=N\), retain only \(j=0\) and odd \(\ell\) in \eqref{eq:R-alias}.  Since
\(M=2^N\), differentiation of \eqref{eq:chi-Phi} at the order-\(N\) zero gives
\[
 \chi^{(N)}(M\ell)=2^{-N}\Phi^{(N)}(2^{N-1}\ell)
 =-N!2^{-N^2}\ell^{-N}\Phi(\ell/2).
\]
Substitution and pairing \(\ell\) with \(-\ell\) give
\eqref{eq:odd-defect-series}.  The sign law \eqref{eq:half-sign} follows by counting the
integer zeros crossed by the real even product \(\Phi\).
\end{proof}

\begin{table}[H]
\centering
\small
\resizebox{0.92\textwidth}{!}{% Generated by dyadic_comb_experiments.py
\begingroup
\renewcommand{\arraystretch}{1.25}
\begin{tabular}{c c r}
\toprule
$N$ & first degree beyond the proved range & exact defect \\
\midrule
1 & 1 & ${-1}/{144}$ \\
2 & 3 & ${-97}/{1382400}$ \\
3 & 3 & ${23}/{22118400}$ \\
4 & 5 & ${9671}/{11985978654720}$ \\
5 & 5 & ${-1543}/{767102633902080}$ \\
6 & 7 & ${-9903527}/{146509414227756711936000}$ \\
7 & 7 & ${1196113}/{37506410042305718255616000}$ \\
8 & 9 & ${21035141003}/{590082134533477495427314445451264000}$ \\
\bottomrule
\end{tabular}
\endgroup
}
\caption{Exact first defects through level eight.  The table supports the exact-precision
conjecture in Conjecture~\ref{conj:exact-degree}.}
\label{tab:first-defects}
\end{table}

\begin{table}[H]
\centering
\small
\resizebox{\textwidth}{!}{% Generated by dyadic_comb_experiments.py
\begin{tabular}{c ccccccccccc}
\toprule
$N\backslash p$ & 0 & 1 & 2 & 3 & 4 & 5 & 6 & 7 & 8 & 9 & 10 \\
\midrule
0 & $\checkmark$ & $\cdot$ & $\cdot$ & $\cdot$ & $\cdot$ & $\cdot$ & $\cdot$ & $\cdot$ & $\cdot$ & $\cdot$ & $\cdot$ \\
1 & $\checkmark$ & $\cdot$ & $\cdot$ & $\cdot$ & $\cdot$ & $\cdot$ & $\cdot$ & $\cdot$ & $\cdot$ & $\cdot$ & $\cdot$ \\
2 & $\checkmark$ & $\checkmark$ & $\checkmark$ & $\cdot$ & $\cdot$ & $\cdot$ & $\cdot$ & $\cdot$ & $\cdot$ & $\cdot$ & $\cdot$ \\
3 & $\checkmark$ & $\checkmark$ & $\checkmark$ & $\cdot$ & $\cdot$ & $\cdot$ & $\cdot$ & $\cdot$ & $\cdot$ & $\cdot$ & $\cdot$ \\
4 & $\checkmark$ & $\checkmark$ & $\checkmark$ & $\checkmark$ & $\checkmark$ & $\cdot$ & $\cdot$ & $\cdot$ & $\cdot$ & $\cdot$ & $\cdot$ \\
5 & $\checkmark$ & $\checkmark$ & $\checkmark$ & $\checkmark$ & $\checkmark$ & $\cdot$ & $\cdot$ & $\cdot$ & $\cdot$ & $\cdot$ & $\cdot$ \\
6 & $\checkmark$ & $\checkmark$ & $\checkmark$ & $\checkmark$ & $\checkmark$ & $\checkmark$ & $\checkmark$ & $\cdot$ & $\cdot$ & $\cdot$ & $\cdot$ \\
7 & $\checkmark$ & $\checkmark$ & $\checkmark$ & $\checkmark$ & $\checkmark$ & $\checkmark$ & $\checkmark$ & $\cdot$ & $\cdot$ & $\cdot$ & $\cdot$ \\
8 & $\checkmark$ & $\checkmark$ & $\checkmark$ & $\checkmark$ & $\checkmark$ & $\checkmark$ & $\checkmark$ & $\checkmark$ & $\checkmark$ & $\cdot$ & $\cdot$ \\
\bottomrule
\end{tabular}
}
\caption{Exact rational check of \eqref{eq:EM-collapse}.  A check mark means that the
complete Euler--Maclaurin polynomial equals the direct dyadic sum exactly.}
\label{tab:exactness-matrix}
\end{table}

\section{Complex, fractional, and negative powers}
\label{sec:complex}

\subsection{Exact Hurwitz-zeta form}

The Bernoulli formula is the natural specialization of a more general finite Hurwitz-zeta
identity.  For \(s\in\C\),
\begin{equation}
 \sum_{k=a+1}^{M-1}k^s
 =\zeta(-s,a+1)-\zeta(-s,M).
 \label{eq:Hurwitz-finite}
\end{equation}
The pole of each individual zeta function cancels in the difference, which is an entire
function of \(s\).

\begin{theorem}[Exact complex-power comb]
\label{thm:Hurwitz-comb}
For every \(\alpha\in\C\),
\begin{align}
 L_{N,\alpha}
 &=C_NM^{-\alpha-1}
 \sum_{a=0}^{M-2}\varepsilon_a
 \sum_{j=0}^{\lfloor N/2\rfloor}\binom N{2j}c_j
 \sum_{r=0}^{N-2j}\binom{N-2j}{r}2^r(-1)^{N-2j-r}
 \notag\\
 &\quad\times
 \sum_{q=0}^{r}\binom rq(-a)^{r-q}
 \left[
 \zeta(-\alpha-q,a+1)-\zeta(-\alpha-q,M)
 \right].
 \label{eq:Hurwitz-comb}
\end{align}
The right side is entire in \(\alpha\), and at \(\alpha=p\in\N_0\) it reduces to
\cref{thm:bernoulli-TM} by
\(
 \zeta(-m,x)=-B_{m+1}(x)/(m+1)
\).
\end{theorem}

\begin{proof}
After exchanging the \(k\)- and \(a\)-sums as in the proof of
\cref{thm:bernoulli-TM}, expand
\[
 (2(k-a)-1)^d
 =\sum_{r=0}^{d}\binom dr2^r(-1)^{d-r}
 \sum_{q=0}^{r}\binom rqk^q(-a)^{r-q}.
\]
The remaining finite sum of \(k^{\alpha+q}\) is \eqref{eq:Hurwitz-finite}.
\end{proof}

\begin{corollary}[Negative integer powers are rational]
\label{cor:negative-rational}
For \(q\in\N\),
\begin{equation}
 \boxed{\displaystyle
 L_{N,-q}=M^{q-1}\sum_{k=1}^{M-1}
 \frac{F(k/M)}{k^q}\in\Q.}
 \label{eq:negative-rational}
\end{equation}
The case \(q=1\) is the harmonic Fabius comb
\(
 \sum_{k=1}^{M-1}F(k/M)/k
\).
\end{corollary}

\subsection{Euler--Maclaurin for fixed complex exponent}

For fixed \(\alpha\), the function \(x^\alpha F(x)\) extends smoothly and flatly through
\(x=0\), because \(F\) and all of its derivatives decay faster than every power there.
At \(x=1\), all positive derivatives of \(F\) vanish, so
\begin{equation}
 \left.\frac{\dd^m}{\dd x^m}\bigl(x^\alpha F(x)\bigr)\right|_{x=1}
 =\fall\alpha m.
 \label{eq:endpoint-derivative-alpha}
\end{equation}
Standard Euler--Maclaurin therefore gives an all-orders expansion, locally uniformly in
\(\alpha\).

\begin{theorem}[Complex-power asymptotic]
\label{thm:complex-EM}
For \(\alpha\) in any compact subset of \(\C\),
\begin{equation}
 \boxed{\displaystyle
 L_{N,\alpha}\sim I(\alpha)-\frac h2+
 \sum_{r\ge1}\frac{B_{2r}}{(2r)!}
 \fall\alpha{2r-1}h^{2r},
 \qquad N\to\infty,}
 \label{eq:complex-EM}
\end{equation}
with a remainder after \(R\) terms that is \(\bigO_K(h^{2R+2})\) on every compact
\(K\subset\C\).  Equivalently, for every prescribed algebraic order \(A\), subtracting
sufficiently many displayed terms leaves \(\bigO_K(h^A)\).
\end{theorem}

\begin{remark}
When \(\alpha=p\in\N_0\), the formal series terminates.  The unusual content of
\cref{thm:EM-collapse} is that the remaining alias remainder is not merely small: it is
identically zero once \(N\) crosses a degree-dependent threshold.
\end{remark}

For \(\alpha=-1\),
\begin{equation}
 L_{N,-1}\sim I(-1)-\frac h2-
 \sum_{r\ge1}\frac{B_{2r}}{2r}h^{2r}.
 \label{eq:alpha-minus-one}
\end{equation}
The table below illustrates convergence for several nonintegral exponents.

\begin{table}[H]
\centering
\small
\resizebox{0.96\textwidth}{!}{% Generated by dyadic_comb_experiments.py
\begin{tabular}{c rrr r}
\toprule
$\alpha$ & $L_{4,\alpha}$ & $L_{6,\alpha}$ & $L_{8,\alpha}$ & EM estimate of $I(\alpha)$ \\
\midrule
-1 & 0.726663799849 & 0.750407155874 & 0.756285603841 & 0.758240000404 \\
-0.5 & 0.575973485504 & 0.599563638374 & 0.605432549961 & 0.607386310743 \\
0.5 & 0.389880928801 & 0.413165845835 & 0.419015684123 & 0.42096817334 \\
1.5 & 0.283877847065 & 0.306857575817 & 0.312688340557 & 0.314639558208 \\
3.5 & 0.17301875663 & 0.195388418551 & 0.201181037431 & 0.203129711955 \\
\bottomrule
\end{tabular}
}
\caption{Direct complex-power combs and a six-term reversed Euler--Maclaurin estimate.
The last column is numerical evidence, not an independent closed-form evaluation.}
\label{tab:fractional-values}
\end{table}

\subsection{A Newton-series alternative}

Let
\begin{equation}
 A_N(z)=\sum_{k=1}^{M-1}F(k/M)z^k.
\end{equation}
Because \(k^\alpha\) has the Newton expansion
\begin{equation}
 k^\alpha=\sum_{r=0}^{k}\binom{k}{r}\Delta^r0^\alpha,
 \label{eq:Newton-power}
\end{equation}
one may write
\begin{equation}
 L_{N,\alpha}=M^{-\alpha-1}
 \sum_{r=0}^{M-1}\Delta^r0^\alpha
 \sum_{k=r}^{M-1}\binom krF(k/M).
 \label{eq:Newton-comb}
\end{equation}
This is finite at every \(N\) and gives a second entire interpolation in \(\alpha\).  It
is potentially advantageous near negative integers, where the Hurwitz-zeta representation
contains large cancelling terms.

\section{Rvachev combs and the critical exponent}
\label{sec:up}

\subsection{Punctured absolute comb}

On \([0,1]\), \(u(x)=1-F(x)\).  Define the pure power comb
\begin{equation}
 P_{N,\alpha}
 =h\sum_{k=1}^{M-1}(k/M)^\alpha
 =M^{-\alpha-1}\bigl[\zeta(-\alpha)-\zeta(-\alpha,M)\bigr].
 \label{eq:pure-comb}
\end{equation}
The punctured symmetric absolute Rvachev comb is
\begin{equation}
 U_{N,\alpha}^{\times}
 :=h\sum_{0<|k|<M}|k/M|^\alpha u(k/M).
 \label{eq:U-punctured}
\end{equation}

\begin{proposition}[Fabius--Rvachev comb relation]
\label{prop:U-F}
For every \(\alpha\in\C\),
\begin{equation}
 \boxed{\displaystyle
 U_{N,\alpha}^{\times}=2\bigl(P_{N,\alpha}-L_{N,\alpha}\bigr).}
 \label{eq:U-relation}
\end{equation}
If the principal branch is used for \(x^\alpha\) on the negative axis, then the signed
complex monomial comb is
\begin{equation}
 h\sum_{0<|k|<M}(k/M)^\alpha u(k/M)
 =(1+\e^{\pi\ii\alpha})\bigl(P_{N,\alpha}-L_{N,\alpha}\bigr).
 \label{eq:U-branch}
\end{equation}
\end{proposition}

\begin{corollary}[Recovery of exact self-sampling moments]
\label{cor:recover-self-sampling}
For \(r\ge1\) and \(2r\le N\),
\begin{equation}
 U_{N,2r}^{\times}=c_r.
 \label{eq:U-even-exact}
\end{equation}
For \(r=0\), adding the central sample \(hu(0)=h\) gives
\begin{equation}
 h+U_{N,0}^{\times}=1.
 \label{eq:U-normalization}
\end{equation}
\end{corollary}

\begin{proof}
This is the centered phase of the corpus-known self-sampling quadrature for \(u\).  It
also follows directly by subtracting the exact Faulhaber sum from
\eqref{eq:EM-collapse}; the Bernoulli boundary polynomials cancel.
\end{proof}

\begin{resultbox}[Consistency bridge]
The corpus's self-weighted quadrature and the present unweighted Fabius comb are distinct
constructions, but \(u=1-F\) makes them complementary.  Exact up-moments are equivalent
to cancellation between an ordinary Faulhaber sum and the exact Euler--Maclaurin
polynomial for \(L_{N,2r}\).
\end{resultbox}

\subsection{Continuous Rvachev Mellin transform}

For \(\Re\alpha>-1\),
\begin{equation}
 J(\alpha):=\int_0^1x^\alpha u(x)\,\dd x
 =\frac1{\alpha+1}-I(\alpha).
 \label{eq:Jalpha}
\end{equation}
Since \(u(x)=1+\bigO(x^A)\) for every \(A>0\) as \(x\downarrow0\), \(J\) extends
meromorphically to \(\C\) with exactly one pole:
\begin{equation}
 \Res_{\alpha=-1}J(\alpha)=1.
 \label{eq:J-residue}
\end{equation}
The full absolute moment is \(2J(\alpha)\).

\subsection{Critical and supercritical negative powers}

At \(\alpha=-1\), \eqref{eq:U-relation} becomes
\begin{equation}
 U_{N,-1}^{\times}=2\bigl(H_{M-1}-L_{N,-1}\bigr),
 \label{eq:U-critical-exact}
\end{equation}
where \(H_n\) is the harmonic number.  The Euler--Maclaurin expansions of
\(H_{M-1}\) and \(L_{N,-1}\) have identical algebraic tails, hence cancel term by term.

\begin{theorem}[Critical logarithm]
\label{thm:critical-log}
As \(N\to\infty\),
\begin{equation}
 \boxed{\displaystyle
 U_{N,-1}^{\times}
 =2\log M+2\bigl(\gamma-I(-1)\bigr)+\bigO(M^{-A})}
 \label{eq:critical-log}
\end{equation}
for every fixed \(A>0\).  Thus the only divergence is the universal logarithm generated
by \(u(0)=1\); the Fabius correction contributes the finite constant \(-2I(-1)\).
\end{theorem}

\begin{proof}
Use
\[
 H_{M-1}=\log M+\gamma-\frac h2-
 \sum_{r=1}^{R}\frac{B_{2r}}{2r}h^{2r}+\bigO(h^{2R+2})
\]
and compare with \eqref{eq:alpha-minus-one}.  Since \(R\) is arbitrary, all algebraic
terms cancel.
\end{proof}

For an integer \(q>1\), generalized harmonic asymptotics similarly give
\begin{equation}
 \boxed{\displaystyle
 U_{N,-q}^{\times}
 =2\zeta(q)M^{q-1}
 -2\left(\frac1{q-1}+I(-q)\right)
 +\bigO(M^{-A})}
 \label{eq:U-q}
\end{equation}
for every fixed \(A>0\).  Again all algebraic Euler--Maclaurin corrections beyond the
finite-part constant cancel.  This is a discrete Mellin-renormalization phenomenon:
\(2\zeta(q)M^{q-1}\) is the lattice singularity, while
\(-2((q-1)^{-1}+I(-q))\) is the finite part of the Rvachev integral.

\section{Iterated sums of arbitrary integer order}
\label{sec:iterated}

\subsection{Definition and binomial kernel}

For a sequence \(a=(a_k)_{k\ge0}\), let
\begin{equation}
 (\Sigma a)_m=\sum_{k=0}^{m}a_k,
 \qquad
 \Sigma^n=\underbrace{\Sigma\circ\cdots\circ\Sigma}_{n\text{ times}}.
 \label{eq:Sigma}
\end{equation}
Take
\begin{equation}
 a_k^{(N,\alpha)}=(k/M)^\alpha F(k/M),
 \qquad 0\le k<M,
 \label{eq:ak}
\end{equation}
with the punctured convention at \(k=0\).  Define the normalized \(n\)-fold comb by
\begin{equation}
 L_{N,\alpha}^{[n]}
 :=h^n(\Sigma^na^{(N,\alpha)})_{M-1}.
 \label{eq:Ln-iter-def}
\end{equation}
Repeated prefix summation gives
\begin{equation}
 \boxed{\displaystyle
 L_{N,\alpha}^{[n]}
 =h^n\sum_{k=1}^{M-1}
 \binom{M-k+n-2}{n-1}(k/M)^\alpha F(k/M).}
 \label{eq:iter-binomial}
\end{equation}
Equivalently, if
\(A_{N,\alpha}(z)=\sum_{k=1}^{M-1}(k/M)^\alpha F(k/M)z^k\), then
\begin{equation}
 (\Sigma^na)_m=[z^m]\frac{A_{N,\alpha}(z)}{(1-z)^n}.
 \label{eq:iter-gf}
\end{equation}
For \(\alpha=p\in\N_0\), one may take
\(A_{N,p}(z)=M^{-p}\Dop^pA_N(z)\).

\subsection{Reduction to shifted one-fold combs}

Use
\begin{equation}
 \binom{M-k+n-2}{n-1}
 =\frac1{(n-1)!}\rise{(M-k)}{n-1}
 =\frac1{(n-1)!}
 \sum_{j=0}^{n-1}\stirlingfirst{n-1}{j}(M-k)^j,
 \label{eq:rising-Stirling}
\end{equation}
where \(\stirlingfirst{r}{j}\) is an unsigned Stirling number of the first kind.

\begin{theorem}[Exact \(n\)-fold reduction]
\label{thm:iterated-reduction}
For \(n\ge1\) and every \(\alpha\in\C\),
\begin{equation}
 \boxed{\displaystyle
 L_{N,\alpha}^{[n]}
 =\frac{M^{1-n}}{(n-1)!}
 \sum_{j=0}^{n-1}\stirlingfirst{n-1}{j}M^j
 \sum_{\ell=0}^{j}(-1)^\ell\binom j\ell
 L_{N,\alpha+\ell}.}
 \label{eq:iterated-reduction}
\end{equation}
Thus an \(n\)-fold sum requires only the \(n\) shifted base combs
\(L_{N,\alpha},\ldots,L_{N,\alpha+n-1}\).
\end{theorem}

\begin{proof}
Expand \((M-k)^j\) in \eqref{eq:rising-Stirling}.  Since
\begin{equation}
 \sum_{k=1}^{M-1}k^\ell(k/M)^\alpha F(k/M)
 =M^{\ell+1}L_{N,\alpha+\ell},
 \label{eq:shift-base}
\end{equation}
substitution into \eqref{eq:iter-binomial} gives the stated powers of \(M\).
\end{proof}

\begin{corollary}[Closed form for natural powers]
\label{cor:iterated-natural}
For \(p\in\N_0\), substitute either \cref{thm:bernoulli-TM} or
\cref{thm:EM-remainder} into \eqref{eq:iterated-reduction}.  In particular, if
\begin{equation}
 N\ge2\left\lceil\frac{p+n-1}{2}\right\rceil,
 \label{eq:iterated-threshold}
\end{equation}
then every shifted comb in \eqref{eq:iterated-reduction} is given by the exact finite
Euler--Maclaurin polynomial \eqref{eq:EM-collapse}.  Hence
\(L_{N,p}^{[n]}\) is an explicit rational polynomial in \(h\), with coefficients formed
from \(I_p,\ldots,I_{p+n-1}\), Bernoulli numbers, and Stirling numbers.
\end{corollary}

\subsection{Continuous limit}

From \eqref{eq:iter-binomial},
\begin{equation}
 h^{n-1}\binom{M-k+n-2}{n-1}
 \longrightarrow\frac{(1-x)^{n-1}}{(n-1)!},
 \qquad x=k/M.
\end{equation}
Therefore
\begin{equation}
 \boxed{\displaystyle
 \lim_{N\to\infty}L_{N,\alpha}^{[n]}
 =\frac1{(n-1)!}\int_0^1(1-x)^{n-1}x^\alpha F(x)\,\dd x}
 \label{eq:iterated-limit}
\end{equation}
for every \(\alpha\in\C\), and binomial expansion gives
\begin{equation}
 \lim_{N\to\infty}L_{N,\alpha}^{[n]}
 =\frac1{(n-1)!}
 \sum_{\ell=0}^{n-1}(-1)^\ell\binom{n-1}{\ell}I(\alpha+\ell).
 \label{eq:iterated-limit-I}
\end{equation}
This is the ordinary Volterra kernel of order \(n\) evaluated at the right endpoint.

\section{Fractional-order summation}
\label{sec:fractional-sum}

The generating-function identity \((1-z)^{-n}\) has a canonical positive-real
continuation.  For \(\rho\in\C\) with \(\Re\rho>0\), define
\begin{equation}
 L_{N,\alpha}^{[\rho]}
 :=h^\rho\sum_{k=1}^{M-1}
 \frac{\Gamma(M-k+\rho-1)}{\Gamma(\rho)\Gamma(M-k)}
 (k/M)^\alpha F(k/M).
 \label{eq:fractional-sum}
\end{equation}
The coefficient is
\([
 z^{M-1-k}](1-z)^{-\rho}
\), so
\begin{equation}
 \boxed{\displaystyle
 L_{N,\alpha}^{[\rho]}
 =h^\rho[z^{M-1}](1-z)^{-\rho}A_{N,\alpha}(z).}
 \label{eq:fractional-coeff}
\end{equation}
At \(\rho=n\in\N\), this is exactly \eqref{eq:iter-binomial}.

\begin{proposition}[Discrete semigroup]
\label{prop:discrete-semigroup}
Let \(\Sigma^\rho\) denote convolution with the coefficients of
\((1-z)^{-\rho}\).  Wherever the finite sums are defined,
\begin{equation}
 \Sigma^\rho\Sigma^\sigma=\Sigma^{\rho+\sigma}.
 \label{eq:discrete-semigroup}
\end{equation}
\end{proposition}

\begin{proof}
Multiply generating functions:
\((1-z)^{-\rho}(1-z)^{-\sigma}=(1-z)^{-\rho-\sigma}\).
\end{proof}

\begin{theorem}[Fractional Volterra limit]
\label{thm:fractional-limit}
For \(\Re\rho>0\) and \(\alpha\) in compact subsets of \(\C\),
\begin{equation}
 \boxed{\displaystyle
 \lim_{N\to\infty}L_{N,\alpha}^{[\rho]}
 =\frac1{\Gamma(\rho)}
 \int_0^1(1-x)^{\rho-1}x^\alpha F(x)\,\dd x.}
 \label{eq:fractional-limit}
\end{equation}
\end{theorem}

\begin{proofidea}
Uniform Stirling asymptotics for the Gamma ratio give
\[
 h^{\rho-1}
 \frac{\Gamma(M-k+\rho-1)}{\Gamma(\rho)\Gamma(M-k)}
 \to\frac{(1-x)^{\rho-1}}{\Gamma(\rho)}.
\]
Endpoint flatness at \(x=0\), together with \(\Re\rho>0\) at \(x=1\), supplies a
uniform integrable majorant after splitting off finitely many boundary nodes.
\end{proofidea}

\begin{remark}
For nonintegral \(\rho\), there is no finite Stirling reduction analogous to
\eqref{eq:iterated-reduction}.  One promising replacement is a convergent Newton series
in shifted combs \(L_{N,\alpha+j}\), with coefficients obtained from the generalized
binomial kernel.  This is developed as an open direction in \cref{sec:frontiers}.
\end{remark}

\section{Arithmetic consequences}
\label{sec:arithmetic}

The corpus supplies explicit common denominator bounds for all level-\(N\) dyadic values.
One convenient nonsharp bound is
\begin{equation}
 \Delta_N
 =N!\,2^{N(N+1)/2}
 (2\lfloor N/2\rfloor+1)!!
 \prod_{j=1}^{\lfloor N/2\rfloor}(2^{2j}-1),
 \label{eq:DeltaN}
\end{equation}
for which
\(
 \Delta_NF(k/M)\in\Z
\) at every level-\(N\) node.

\begin{proposition}[Immediate denominator bounds]
\label{prop:denominator-bounds}
For \(p\in\N_0\),
\begin{equation}
 \Delta_NM^{p+1}L_{N,p}\in\Z.
 \label{eq:denom-positive}
\end{equation}
For \(q\in\N\), if
\(
 \Lambda_M=\operatorname{lcm}(1,\ldots,M-1)
\), then
\begin{equation}
 \Delta_N\Lambda_M^qL_{N,-q}\in\Z.
 \label{eq:denom-negative}
\end{equation}
More generally, the same bounds propagate through
\eqref{eq:iterated-reduction}, with the additional factorial and powers of \(M\) visible
there.
\end{proposition}

\begin{remark}
These bounds are deliberately crude.  The exact Euler--Maclaurin formula often produces
large cancellations between the moment denominator in \(I_p\) and the Bernoulli powers
of \(2\).  The sharp \(2\)-adic and odd-denominator valuations appear to be governed by
an interaction among Mersenne factors \(2^{2j}-1\), Bernoulli denominators, and the
Prouhet block.  See Problem~\ref{prob:sharp-valuations}.
\end{remark}

\section{Numerical and symbolic verification}
\label{sec:experiments}

The companion program \path{dyadic_comb_experiments.py} uses three independent exact
routes:
\begin{enumerate}
\item direct dyadic samples from \eqref{eq:dyadic-convolution};
\item the exchanged Bernoulli--Thue--Morse formula \eqref{eq:L-H};
\item the Bell--Prouhet coefficient collapse \eqref{eq:Bell-collapse}.
\end{enumerate}
It also constructs the rational OGF \eqref{eq:Frow} symbolically, checks the exact
Euler--Maclaurin range through level eight, compares direct and reduced iterated sums,
and evaluates nonintegral exponents at high precision.

\begin{statusbox}[Reproducibility]
Running
\begin{lstlisting}[style=python]
python dyadic_comb_experiments.py --outdir generated --max-level 8
\end{lstlisting}
performs all self-checks and regenerates the CSV and LaTeX tables included in this paper.
All exact checks use \texttt{fractions.Fraction}; SymPy is used for Bernoulli polynomials,
Bell coefficient extraction, and rational-function simplification; mpmath is used only
for nonintegral powers.
\end{statusbox}

The exact checks performed by the script include
\begin{equation}
 L_{N,p}^{\mathrm{direct}}
 =L_{N,p}^{\mathrm{Bernoulli}}
 =L_{N,p}^{\mathrm{Bell}}
 \label{eq:three-way-check}
\end{equation}
for representative ranges, and
\begin{equation}
 L_{N,p}^{[n],\mathrm{direct}}
 =L_{N,p}^{[n],\mathrm{Stirling}}
 \label{eq:iterated-check}
\end{equation}
for \(N\le6\), \(p\le4\), and \(n\le4\).  No floating-point comparison is used to
assert an exact theorem.

\section{Conjectures and frontier directions}
\label{sec:frontiers}

\subsection{Exact degree of Euler--Maclaurin collapse}

The proved range is \(p\le d_N=2\lfloor N/2\rfloor\).  Exact computations through
\(N=8\) show a nonzero defect at the first possible degree beyond it.

\begin{conjecture}[Exact precision]\label{conj:exact-degree}
For every \(N\ge1\),
\begin{equation}
 \defect_{N,p}=0\quad(0\le p\le d_N),
 \qquad
 \defect_{N,d_N+1}\ne0.
 \label{eq:exact-degree-conj}
\end{equation}
Equivalently, the Euler--Maclaurin collapse has exact maximal degree
\(2\lfloor N/2\rfloor\).
\end{conjecture}

For odd \(N\), this reduces to nonvanishing of the explicit series
\eqref{eq:odd-defect-series}.  The first term is positive in magnitude and the Fourier
product decays rapidly, suggesting a first-harmonic domination argument.  A rigorous proof
must control possible Thue--Morse-signed cancellation uniformly in \(N\).  For even
\(N\), the first defect mixes the first two odd-alias jets with the leading twice-odd
alias layer; the local Bell-polynomial zero-jet calculus from the corpus should make this
finite-layer structure explicit.

\subsection{Sign and asymptotic size of the first defect}

\begin{conjecture}[Odd-level sign]
For odd \(N\), the sum in \eqref{eq:odd-defect-series} is positive.  Hence
\begin{equation}
 \sgn\defect_{N,N}=(-1)^{(N+1)/2}.
 \label{eq:odd-sign-conj}
\end{equation}
\end{conjecture}

A stronger target is an asymptotic expansion dominated by \(\Phi(1/2)\):
\begin{equation}
 \defect_{N,N}
 \sim\frac{2(-1)^{(N+1)/2}N!\Phi(1/2)}
 {(2\pi)^{N+1}2^{N(N+1)}}.
 \label{eq:defect-first-harmonic}
\end{equation}
The correction is exponentially small in \(N\) from the odd denominators and is further
suppressed by the sinc-product decay.  The repository's sharp arithmetic-ray estimates
should be adaptable to a quantitative proof.

\subsection{Row-polynomial geometry}

Write
\begin{equation}
 \Frow_N(z)=\frac{Q_N(z)}{1-z},
 \qquad \deg Q_N\le M.
 \label{eq:QN-row}
\end{equation}
The first examples exhibit reciprocal factors, powers of \(1+z\), and unexpectedly
small-coefficient palindromic pieces.  Natural problems are:
\begin{problem}
Determine the exact reciprocal symmetry of \(Q_N\), the multiplicities of cyclotomic
factors, and whether the noncyclotomic zeros obey a limiting curve or annulus law.
\end{problem}
A route through \(F(k/M)+F((M-k)/M)=1\) gives a functional equation for \(A_N(z)\), while
\eqref{eq:Frow} expresses \(Q_N\) through the Thue--Morse mask and the moment polynomial
\(P_N(2\Dop-1)\).  The collision of those two descriptions may lead to a finite spectral
factorization.

\subsection{Sharp denominator valuations}

\begin{problem}
\label{prob:sharp-valuations}
Determine
\begin{equation}
 \vtwo\!\left(\operatorname{den}L_{N,p}\right)
 \quad\text{and}\quad
 v_\ell\!\left(\operatorname{den}L_{N,p}\right)
 \label{eq:valuation-problem}
\end{equation}
for odd primes \(\ell\), uniformly in \(N,p\).  Do the same for the iterated sums and
negative powers.
\end{problem}
The exact Euler--Maclaurin formula reduces the stable range to the arithmetic of
\(I_p\), Bernoulli denominators, and powers of two.  Outside that range, the Bell--Prouhet
formula adds only \(p+2\) signed moments and may expose a tractable valuation filtration.

\subsection{Double scaling near negative powers}

For fixed \(q\), \cref{sec:up} gives the complete leading behavior.  A different regime
arises when \(q=q_N\) grows with \(N\).  The smallest lattice node contributes
\begin{equation}
 M^{q_N-1}F(1/M),
 \label{eq:smallest-node}
\end{equation}
and \(F(2^{-N})\) has a log-squared/Lambert-\(W_{-1}\) asymptotic.  Balancing that term
against the next nodes should produce a phase diagram in the ratio \(q_N/N\), with
transitions governed by a discrete saddle over the lattice index.  This is a natural
meeting point of the repository's endpoint asymptotics and the present comb calculus.

\subsection{Fractional summation as a two-variable special function}

The object
\begin{equation}
 (\rho,\alpha)\longmapsto L_{N,\alpha}^{[\rho]}
 \label{eq:two-var-special}
\end{equation}
is entire in \(\alpha\) and meromorphic in \(\rho\) before cancellation.  Promising
directions include:
\begin{itemize}
\item a Newton series in shifted \(L_{N,\alpha+j}\) for nonintegral \(\rho\);
\item analytic continuation in \(\rho\) after subtracting the first lattice jets;
\item a Mellin--Barnes representation for the Gamma kernel;
\item a uniform double limit \(N\to\infty\), \(\rho=\rho_N\), including large-order
      fractional Volterra asymptotics;
\item exact arithmetic at rational \(\rho\), where Gamma values may organize into
      controlled transcendental factors times rational dyadic data.
\end{itemize}

\subsection{\texorpdfstring{Base-\(b\)}{Base-b} and digital generalizations}

The Bell--Prouhet mechanism depends only on a finite digital product and a polynomial
kernel.  Replace
\(
 \prod_{j<N}(1-z^{2^j})
\)
by a base-\(b\) digit mask or a root-of-unity weighted product.  Whenever the mask has a
zero of order \(N\) at \(z=1\), the same argument yields a short Bell collapse.  The
corresponding probability law would be a geometrically weighted sum of uniform or
discrete digits.  One may ask which masks also create high-order zeros at the dual
lattice, and hence an exact Euler--Maclaurin range.

\subsection{Formalization roadmap}

The most modular Lean targets are finite and algebraic:
\begin{enumerate}
\item formalize the convolution OGF \eqref{eq:Frow} and the pole-cancellation degree
      bound;
\item formalize \eqref{eq:theta-Bell-egf} as a finite-product identity and derive
      \eqref{eq:TN-Bell};
\item formalize the Bernoulli sum exchange \eqref{eq:L-H};
\item build the periodic trapezoidal alias theorem for a function with one endpoint jump;
\item combine it with the already formalized integer zero orders of \(\Phi\) to prove
      \cref{thm:EM-collapse};
\item formalize the finite Stirling reduction \eqref{eq:iterated-reduction};
\item only then add the analytic Hurwitz and fractional-order layers.
\end{enumerate}
The exactness theorem is particularly suitable for formalization because its essential
spectral argument is finite after the zero-order lemma: every forbidden alias derivative
is literally zero.

\section{Conclusion}

The ordinary dyadic Riemann sums of \(x^pF(x)\) hide considerably more structure than a
generic discretization error.  Their exact dyadic arithmetic produces a rational
sample-row generating function; Prouhet cancellation compresses exponentially many signs
to finitely many Bell polynomials; and the same dyadic geometry forces an exact
Euler--Maclaurin formula on a degree range that grows linearly with the level.  The first
failure is not an opaque remainder but a controlled sinc-product Dirichlet series, and
negative or fractional powers remain accessible through a finite Hurwitz-zeta formula.
Iterated summation adds no essentially new arithmetic obstruction: unsigned Stirling
numbers reduce it to a short list of shifted base combs.

The strongest immediate next step is to prove exact maximality of the degree range.  At
odd levels this is a concrete positivity/noncancellation problem for the half-integer
Fourier product; at even levels it is a finite competition among the first two 2-adic
alias layers.  Both are tightly connected to machinery already developed in the
repository, making them realistic frontier targets rather than vague extrapolations.

\appendix

\section{Derivation of the row symmetry}
\label{app:row-symmetry}

Let
\begin{equation}
 A_N(z)=\sum_{k=0}^{M-1}F(k/M)z^k.
\end{equation}
Using \(F((M-k)/M)=1-F(k/M)\), one obtains
\begin{align}
 z^MA_N(z^{-1})
 &=\sum_{k=1}^{M}F((M-k)/M)z^k
 \notag\\
 &=\sum_{k=1}^{M}\bigl(1-F(k/M)\bigr)z^k
 \notag\\
 &=\frac{z(1-z^M)}{1-z}-A_N(z)-z^M.
 \label{eq:AN-reciprocal}
\end{align}
Equivalently,
\begin{equation}
 A_N(z)+z^MA_N(z^{-1})
 =\frac{z(1-z^{M-1})}{1-z}.
 \label{eq:AN-reciprocal-clean}
\end{equation}
This exact reciprocal relation explains the palindromic features visible in the low-level
row numerators.

\section{First few moments and continuous integrals}
\label{app:moments}

The recurrence \eqref{eq:c-recurrence} begins
\begin{equation}
 c_0=1,
 \qquad c_1=\frac19,
 \qquad c_2=\frac{19}{675},
 \qquad c_3=\frac{583}{59535},
 \qquad c_4=\frac{132809}{32531625}.
 \label{eq:first-c}
\end{equation}
From \eqref{eq:dm-from-c},
\begin{equation}
 d_1=\frac12,
 \quad d_2=\frac5{18},
 \quad d_3=\frac16,
 \quad d_4=\frac{143}{1350},
 \quad d_5=\frac{19}{270}.
 \label{eq:first-d}
\end{equation}
Consequently \eqref{eq:Ip} yields the constants displayed in
\cref{tab:integer-closed}.

\section{Algorithmic summary}
\label{app:algorithm}

For exact computation of \(L_{N,p}\):
\begin{enumerate}
\item compute \(c_0,\ldots,c_{\lfloor N/2\rfloor}\) from
      \eqref{eq:c-recurrence};
\item choose one of three routes:
  \begin{enumerate}[label=(\alph*)]
  \item evaluate all dyadic samples by \eqref{eq:dyadic-convolution} and sum directly;
  \item evaluate \eqref{eq:L-H};
  \item form \(H_{N,p}(a)\), retain only coefficients of degrees \(N\) through
        \(N+p+1\), and apply \eqref{eq:TN-Bell};
  \end{enumerate}
\item when \(p\le2\lfloor N/2\rfloor\), bypass all dyadic samples and use the exact
      Euler--Maclaurin polynomial \eqref{eq:EM-collapse};
\item for an \(n\)-fold sum, first compute
      \(L_{N,p},\ldots,L_{N,p+n-1}\), then apply
      \eqref{eq:iterated-reduction}.
\end{enumerate}
The fourth route is asymptotically constant-time in \(N\) at fixed \(p\), apart from the
cost of the moments entering \(I_p\).

\section{Source and literature notes}
\label{app:sources}

\begin{thebibliography}{99}

\bibitem{Arias1982}
J.~Arias de Reyna,
\newblock \emph{An infinitely differentiable function with compact support: definition
and properties},
\newblock English translation of the 1982 paper, arXiv:1702.05442, 2017.

\bibitem{AriasArithmetic}
J.~Arias de Reyna,
\newblock \emph{Arithmetic of the Fabius function},
\newblock \emph{INTEGERS} \textbf{18} (2018), Article A51; arXiv:1702.06487.

\bibitem{AlloucheShallit}
J.-P.~Allouche and J.~Shallit,
\newblock \emph{Automatic Sequences: Theory, Applications, Generalizations},
\newblock Cambridge University Press, 2003.

\bibitem{Comtet}
L.~Comtet,
\newblock \emph{Advanced Combinatorics: The Art of Finite and Infinite Expansions},
\newblock Reidel, 1974.

\bibitem{NIST}
F.~W.~J. Olver et al. (eds.),
\newblock \emph{NIST Digital Library of Mathematical Functions},
\newblock chapters on Bernoulli polynomials, Hurwitz zeta, and asymptotic expansions.

\bibitem{BerlineVergne}
N.~Berline and M.~Vergne,
\newblock Local asymptotic Euler--Maclaurin expansion for Riemann sums over a
semi-rational polyhedron,
\newblock arXiv:1502.01671, 2015.

\bibitem{BuchheitKessler}
A.~A. Buchheit and T.~Ke\ss ler,
\newblock Singular Euler--Maclaurin expansion,
\newblock arXiv:2003.12422, 2020.

\bibitem{ProveItPrimary}
V.~Reshetnikov,
\newblock \emph{The Fabius Function and Rvachev's Up-Function},
\newblock living manuscript in \repo{}, repository snapshot 28 August 2026.

\bibitem{ProveItSampling}
V.~Reshetnikov,
\newblock \emph{Inverse and Sampling Frontiers for the Fabius--Rvachev System},
\newblock consolidated manuscript in \repo{}, 2026.

\bibitem{ProveItIntegration}
V.~Reshetnikov,
\newblock \emph{Integration and Transform Frontiers for the Fabius--Rvachev System},
\newblock consolidated manuscript in \repo{}, 2026.

\bibitem{ProveItRepresentations}
V.~Reshetnikov,
\newblock \emph{Representation Frontiers for the Fabius--Rvachev System} and
\emph{The Second Wave},
\newblock consolidated manuscripts in \repo{}, 2026.

\end{thebibliography}

\end{document}
